\documentclass[twocolumn]{article}

% --- 基本宏包 ---
\usepackage[UTF8]{ctex}
\usepackage{graphicx}
\usepackage{amsmath}
\usepackage{amsfonts}
\usepackage{amssymb}
\usepackage{booktabs} % 用于高质量表格
\usepackage[colorlinks=true, allcolors=blue]{hyperref} % 用于超链接
\usepackage{pifont}

% --- 页面设置 ---
\usepackage[margin=0.8in]{geometry}

% --- 标题和作者 ---
\title{Trident-Former: Physically-Decoupled Underwater Enhancement Fueled by a High-Fidelity Synthetic Dataset}
\author{Your Name \\ Your Institution \\ \texttt{your.email@example.com}}
\date{\today}

\begin{document}

\maketitle

% --- 摘要 ---
\begin{abstract}
水下图像增强在各类水下视觉任务中扮演着关键角色，其效果直接影响着资源勘测、生态监测与自主导航等系统的环境感知能力与决策可靠性。然而，由于真实水下配对数据的严重匮乏，当前大多数方法难以充分建模水体中光的吸收、散射与折射所导致的复杂退化机制，导致网络在面对严重色偏、结构模糊与纹理丢失等场景时表现不稳定，难以实现真实感与结构保真的统一。

为缓解数据瓶颈并提升网络的物理一致性建模能力，本文构建了一个高保真三维可控水下图像增强数据集。该数据集基于水下光学传播模型与光线追踪渲染技术，生成具备严格物理一致性的成对样本，覆盖清晰图像、高精度深度图以及多退化等级的观测图像，全面模拟水下成像过程中的颜色衰减与空间模糊特性，该数据集在场景丰富度、几何复杂度、物理一致性、退化可控性、标注完备性、图像分辨率以及规模上超越现有所有水下图像增强数据集。

在此基础上，本文进一步提出了一种面向真实感重建的端到端增强框架——Trident-Former。该框架引入双次前向协同机制，构建了双分支物理调制解码器，分别对应两类核心退化：基于 Beer-Lambert 散射模型的颜色校正分支用于恢复光谱失真，基于点扩散函数（PSF）模型的去模糊分支用于重建纹理与边缘细节。两分支通过跨尺度语义融合模块实现信息互补，协同完成色彩与结构的联合优化。



\end{abstract}

% --- 正文 ---
\section{引言}
水下光学成像与单目深度估计在海洋探索、水下机器人导航、海洋生物研究及水下考古等领域扮演着至关重要的角色。然而，与陆地环境相比，水下光照条件极为复杂。
光在水体中传播时会经历强烈的吸收和散射效应：波长较长的红光首先被吸收，导致图像整体偏蓝或偏绿；悬浮颗粒物则引起前向散射（使细节模糊）和后向散射（形成类似面纱的效应），
这些现象共同导致水下图像对比度低、颜色失真、细节丢失，且随着图像深度增加，退化问题会呈指数级加剧。
这些退化不仅严重影响人类的视觉感知，也极大地限制了下游计算机视觉算法（如目标检测、语义分割、三维重建等）在水下环境中的应用效果。

针对水下图像增强问题，研究者们已提出众多方法。早期方法主要基于物理模型，如暗通道先验（DCP）及其变种、Retinex理论等，它们试图通过反演水下成像模型来恢复清晰图像。
尽管这类方法在特定场景下有效，但效果严重依赖模型假设和参数设置，泛化能力有限。近年来，深度学习方法逐渐成为主流，通过数据驱动方式学习退化-恢复映射，实现更为灵活的增强。然而，这些方法仍面临诸多瓶颈：首先，现有公开数据集普遍存在规模不足、场景单一、退化类型单调、几何对齐性差、缺乏高精度物理监督信息等问题，无法为复杂退化机制与高精度结构感知的网络提供充分的训练保障；其次，真实水下环境中退化强度连续变化、类型多样且空间异质，现有模型难以适应这种非线性、高度耦合的退化模式；此外，结构信息与色彩恢复任务存在显著的频率与语义层级差异，二者高度耦合导致网络结构复杂且难以实现整体最优。

为解决上述难题，我们首次构建了一个具备大规模、多退化等级与多模态一致性的高保真水下合成图像增强数据集，在数据信息的广度与深度上均达到了划时代的规模与完备程度，
为复杂结构感知与物理一致性建模的增强网络提供了坚实的数据基础，显著提升了模型训练的覆盖性与泛化能力。
我们基于 Blender 渲染平台，构建并精细设计了多个具有不同水体特性与几何结构的三维水下场景，
在每个场景中针对物体布局、材质属性、水体浑浊度、光照条件与相机路径等关键因素进行了逐项建模与精调控制，模拟了光在水中传播过程中的吸收与散射效应，合成生成了具有真实视觉退化特性的多退化等级序列的图像对，
包含清晰图、不同程度退化图、精准深度图与法线图。
除了服务于图像增强模型的训练，所构建的数据集还同步提供了诸如相机位姿等结构化标注信息，可广泛支持多种几何感知与结构理解任务，包括多视角三维重建、神经辐射场建模、新视角合成、相机姿态回归以及视觉SLAM算法评估等，展现出良好的通用性与跨任务迁移能力。
这一数据集不仅极大丰富了水下图像增强的研究资源，也为推动复杂场景下物理一致性建模和高精度结构感知网络的发展奠定了坚实基础，树立了该领域数据构建的新标杆。
基于所构建的数据集，我们提出了一种集"频率解耦建模、深度引导融合与多退化协同学习"于一体的高适应性水下图像增强网络架构 \textbf{Trident-Former}。
该网络引入频率解耦的双分支结构，针对低频颜色偏移的与高频结构模糊分别建模，将Beer-Lambert水下成像模型及PSF模糊成像模型融合深度先验嵌入不同的解码路径，实现颜色校正与细节恢复的协同优化。
区别于传统仅能应对离散退化场景的方案，所提网络创新性引入多退化输入建模机制，使模型能够在训练阶段学习到多源退化样本分布，显著提升了在复杂真实场景下对未知退化的自适应与迁移能力。
"双次前向机制"下，深度信息不仅作为辅助先验，更作为主任务核心参与者被深度嵌入网络主干，深度预测与增强过程高度耦合，充分发挥物理与语义双重先验的协同优势，全流程无中断及信息损失。
推理阶段，模型仅需输入单张退化图像，即可同时完成颜色还原与结构重建，还可输出高精度深度图像，为其他深度感知任务（如三维重建、语义分割辅深度等）提供有力支持。


\section{相关工作}
\subsection{水下图像增强数据集与合成策略}
真实水下图像的"退化–清晰"配对数据难以获得，因此大多数水下图像增强任务的数据集构建遵循三个主要路径：
一是收集真实图像并提供主观参考图像作为"伪真值"，二是基于物理模型从清晰图像合成退化图像对 ，三是直接构建合成真实感图像生成\cite{ref1, ref2}。

首先第一类方法有如下代表，UIEB \cite{ref3} 是首个大规模水下图像增强基准数据集，包含 950 张真实水下照片，其中 890 张配有人为挑选或调整生成的参考图像。
参考图像通过专业人员从多种增强算法结果中选优或手工修正获得，构成了近似"主观真值"。
EUVP \cite{ref4} 提供了12K张配对与8k张非配对子集：
配对图像通过调整照明或拍摄参数（如闪光、靠近拍摄等）获得"高质量版本"，非配对图像集合支持 CycleGAN 等无监督训练。
Peng 等人[TIP, 2023] 构建了 LSUI 数据集，该数据集包含 4,279 对真实水下图像与参考图像对，并配有每幅图像的 语义分割图 和 介质透射图。
RUIE \cite{ref10} 进一步从东海实地采集了 4,000 余张图像，并划分为 UCIQE 分数排序子集（UIQS）、边缘保留子集（UCCS）与目标可检测性子集（UHTS）。
这类数据集具备真实感强、来源多样的优点，但其参考图像并非真实无退化图，在几何上并非严格对齐。

第二类方法类方法依赖物理建模合成退化图像。UWCNN \cite{ref5} 在 NYU-V2 等室内 RGB-D 数据上应用 Beer–Lambert 光衰减定律与 Jerlov 水体模型,生成了1k+对合成数据集，该方法对红绿蓝通道分别设置衰减系数，以模拟不同类型水体（共 10 种）中的图像退化。
WaterGAN \cite{ref6} 使用生成对抗网络将清晰 RGB-D 图像转化为水下图像，通过学习的散射模型与物理渲染模块组合实现无监督退化学习。
为了提供真实深度数据支持物理反演，Sea-thru \cite{ref7} 采集了1157对真实水下立体图像，并重建每像素深度图，再结合色彩卡与衰减曲线反推清晰图像。


最后是第三类合成数据集方法。
VAROS \cite{ref8} 基于 Blender 引擎构建了一个水下探测场景，首次在渲染环境中提供高度真实的水下图像序列及其参考图像、深度图、法线图与 IMU 数据。

Atlantis \cite{ref9} 则提出一种基于深度条件扩散模型的生成式合成方案，通过在陆地 RGB-D 数据集（如 KITTI 或 NYU）上学习"深度到水下图像"的扩散映射，
实现结构一致性下的外观风格转化。。

\subsection{图像增强方法}
经典水下图像增强方法多基于光学成像模型的反演过程，尝试从图像中恢复清晰图像与散射参数。早期工作受大气图像去雾模型启发，通常将水下成像近似为光强随距离的指数衰减过程，
并附加后向散射项 \cite{ref1}。尽管这些方法在物理一致性方面具有优势，但普遍依赖均匀水体或稳定照明的假设，在实际复杂场景中往往难以保持鲁棒性 \cite{ref8}。近年来，深度神经网络逐渐成为水下图像增强的主流工具，凭借强大的特征表达与端到端学习能力，能够拟合复杂的非线性退化映射。
与传统方法相比，深度模型虽在性能上取得显著提升，但也面临可解释性差与物理一致性不足的问题。
为弥合这一差距，当前主流研究路径主要沿两类方向展开：其一是多分支融合策略，通过集成多版本输入或网络模块以强化对多种退化因素的感知与修复；
其二是先验信息引导机制，引入物理深度、透射图或水体参数等条件信息，以增强网络对不同区域退化特性的适应能力。这两类策略已成为当前水下图像增强方法的重要发展趋势。

\subsubsection{多分支融合策略方法}
Ancuti 等人最早提出将多种增强版本（如白平衡、Retinex、对比度拉伸）融合为最终图像，
融合权重基于饱和度、清晰度与亮度等局部指标生成 \cite{ref16}。Tong 等人提出在输入图像上分别施加 CLAHE 与多尺度 Retinex，构建增强序列并进行加权融合 \cite{ref17}。
Li 等 (2019) 提出的 WaterNet 利用传统算法生成的多版本输入（如白平衡、直方图均衡、伽马校正），由CNN融合输出高质量图像。
Li 等 (2020) 提出 UWCNN，基于水下先验合成大规模训练集并训练卷积网络，对不同水体退化分别建模，提升了各类场景的泛化能力。
Islam 等 (2020)开发了轻量级GAN模型 FUnIE-GAN 实时增强水下图像,其特点是在保证效率的同时改善感知质量。
此外，Cong 等 (2023) 提出的 PUGAN 将物理成像模型嵌入GAN框架：包括参数估计子网用于物理模型反演生成颜色校正图，再与主增强网络交互，辅以内容和风格双判别约束，兼顾物理真实性和视觉效果。

\subsubsection{引入先验信息引导的增强方法}
先验信息引导融合深度图、介质透射图、水体类型或物理模型估计结果等外部条件，以显式控制网络增强路径 \cite{ref16, ref19, ref21}。
IBLA 方法则结合了颜色模糊度与散射理论，尝试通过图像清晰度估计场景深度，从而恢复图像 \cite{ref7}。Sea-thru 方法使用结构光或多视角图像获取每个像素对应的物理深度，并将其代入修正后的成像模型实现颜色反演 \cite{ref18}。
UIE-Net (Li et al., 2021) 引入介质透射图估计，将介质透射先验与多颜色空间嵌入融合，实现物理约束的端到端增强。
UVZ（Underwater Variable Zoom）方法则采用两阶段结构，先使用深度网络估计图像深度，再根据预测深度对图像不同区域执行增强操作 \cite{ref19}。
这种基于深度的区域自适应机制能够提升远近区域的恢复差异性。扩散模型的引入进一步丰富了条件生成建模的可能性。
Osmosis 是首个将 RGBD 扩散先验与物理模型结合的工作，通过对 RGBD 空间中的无条件扩散模型进行训练，在推理阶段，通过物理一致性的重建误差计算图像梯度，引导扩散采样向后验分布迭代收敛。
将条件信息融入扩散采样过程 \cite{ref23, ref24}。Peng 等人 [TIP 2023] 提出了面向水下图像增强任务的 U-Shape Transformer 网络，引入通道级与空间级注意力机制以建模非一致性退化特性。


\subsection{水下深度估计算法}
由于水下图像的深度配对数据更加难以获得，因此大多数水下图像深度估计会使用一些泛化能力较强的空中深度估计模型，例如ＭｉＤａＳ，MiDaS对未知场景给出合理相对深度，但在水下会因颜色偏差导致过度估计远处深度
（模糊区域深度易受散射影响）。ＵＷ－Ｎｅｔ是首个无监督单张水下深度ＣＮＮ模型，通过ＧＡＮ生成仿真水下图像来训练深度估计。ＵＷ－ＧＡＮ将物理成像模型融入判别器约束。同时输出增强图像和深度图。
Amitai等人(ICRA 2023)提出利用视频帧匹配进行自监督水下深度学习；Wenjie Hu等(Neurocomputing 2022)采用无真值的对偶网络实现水下自监督深度估计。
这类方法在无标签下取得了一定精度。
UDepth (Boxiao Yu等, ICRA 2023)专为水下机器人导航设计的实时单目深度网络。采用轻量Encoder-Decoder并融合水下成像物理约束，实现了快速且稳健的深度推断。UDepth在SQUID数据集上相比空中模型取得更优结果。
Ye等人在TCSVT 2023提出的立体适配网络（利用双目数据蒸馏单目深度）、TreibitZ组的Sea-thru系列利用真实距离约束网络（如Sea-thru-NeRF, CVPR 2023)、以及2024年提出的PUDE物理引导迁移学习方法等。这些方法充分利用水下成像物理和少量带深度数据进行域适配，显著缩小了水下与空中模型性能差距。
测试评估通常在Sea-thru (D3/D5)和SQUID数据集上进行，这两套真实数据已成为水下深度估计的新基准。

\section{Proposed Dataset}

为了支持物理一致性建模与跨退化泛化学习，我们提出了一个高保真、可控性强、具备多模态监督的合成水下图像数据集。

首先，我们构建了多个具有不同几何结构与水体属性的三维水下场景，涵盖自然地形、人造物体与多尺度纹理组合。
然后，我们使用基于物理渲染的光传播建模策略，模拟水体中的吸收、散射、浑浊度与动态光照等关键退化因素，并通过调控水体参数与相机轨迹，实现从轻微到极端退化的多等级控制。
最后，我们对每一帧图像渲染了清晰图、多个退化版本图、深度图与法线图，确保跨模态对齐和高精度几何一致性。

我们的数据集采用严格命名规则与标准化格式组织，支持单目/多视角配置与高频轨迹采样,采集过程中严格控制了相机轨迹的空间覆盖范围，并导出整理了相机内外参，可直接用于多模态监督学习与跨任务模型评估。该数据构建流程为后续水下图像增强与三维重建等任务提供了系统性、高精度、可扩展的数据基础。

\subsection{场景建模与水体环境构建}

在场景建模方面，我们系统性构建了覆盖多维属性的典型水下环境，确保数据集具备高度的现实感与任务通用性。三维场景涵盖淡水与海水场景，分别模拟湖泊、内湾、近岸与热带开放海域；空间结构上包含浅水、深水、洞穴与开放区域，光照条件则从全自然光照（如浅滩日照）过渡到完全依赖人造照明的低照度环境（如洞穴、深水区）。

\begin{figure*}[htbp]
    \centering
    \includegraphics[width=0.95\linewidth]{figs/scene_diversity_showcase.pdf}
    \caption{(a)所构建的典型任务场景多样性示例，包括沉船、桩基、珊瑚礁、海底遗迹等。(b)材质贴图细节与高光通道可视化，展现真实光照响应效果。(c)从24mm到85mm焦段之间多个经典焦段连续视角对比}
    \label{fig:scene_diversity}
\end{figure*}

% \begin{figure}[htbp]
%     \centering
%     \includegraphics[width=0.95\linewidth]{figs/texture_quality.pdf}
%     \caption{材质贴图细节与高光通道可视化，展现真实光照响应效果。}
%     \label{fig:texture_quality}
% \end{figure}

% \begin{figure}[htbp]
%     \centering
%     \includegraphics[width=0.95\linewidth]{figs/24mm-85mm.pdf}
%     \caption{从24mm到85mm焦段之间多个经典焦段连续视角对比}
%     \label{fig:texture_quality}
% \end{figure}


如图~\ref{fig:scene_diversity} 所示，构造来源结合自然生态与人工设施，包括沙地、沉木、珊瑚礁、沉船、港口桩基与管道结构等，任务语义覆盖工业巡检、水产养殖、水下考古与生态监测等实际应用情境。


几何结构方面，我们通过程序化建模融合地形地貌与高频人工构件，构建出层次丰富、遮挡关系明确的多尺度空间布局。自然区域包含岩壁、沙洲、坡面与海沟等典型地貌，呈现连续性地形变化与非规则边界特征；人工构造则包含管道系统、钢结构、残骸、螺栓与浮标等组件，具有明确几何规律与纹理密度。

为进一步增强真实感与视觉多样性，我们引入了植物、水草、沉积物与漂浮颗粒等自然干扰元素，构建动态多扰动环境，模拟真实水体中的光衰减与浑浊动态。材质与纹理设计方面，我们将岩石、沙地、藻类等高分辨率纹理贴图合理铺设于不同地形区域，确保材质分布非重复、光学属性一致。所有纹理均采用 $4\mathrm{K}$ 分辨率，并支持多通道的粗糙度、反射与高光设计，以适应真实渲染下微观光照的动态变化。如图~\ref{fig:texture_quality} 所示，材质表面的微观反射细节在真实渲染中得以高度还原，有效增强图像的纹理层次与真实感。此外，在视角设计上，我们结合常见移动设备（如运动相机与手机）所具备的全画幅 $24\mathrm{mm}$ 广角视角与经典相机的 $85\mathrm{mm}$ 特写视角，覆盖近距离观察与大场景感知两类典型任务视角，进一步增强了场景构建的实用性与沉浸感。



通过以上多维建模策略的协同设计，我们构建了一个具备广泛视觉分布、结构复杂性与物理一致性的真实感水下合成场景集合，为多任务增强模型提供了可靠、高保真、可扩展的数据基础。

\begin{table*}[htbp]
\centering
\caption{各水下图像增强数据集关键特征，对比主流数据集的特征覆盖情况，并以对号（\ding{51}）与叉号（\ding{55}）对比是否具备特定维度属性。维度包括图像是否具备高分辨率（High Res.）、是否提供多模态标注信息（Multimodal Labels）、是否包含相机位姿信息（Camera Pose）、是否为结构对齐的成对图像（Paired Images）、是否覆盖多等级退化类型（Multi Degrad.），以及退化过程是否符合水下成像的物理规律（Physically Plausible）。}
\label{tab:dataset_comparison_vertical}
\begin{tabular}{|l|c|c|c|c|c|c|c|c|}
\hline
\textbf{Feature} & \textbf{UIEB} & \textbf{EUVP} & \textbf{LSUI} & \textbf{UWCNN} & \textbf{Sea-thru} & \textbf{VAROS} & \textbf{RUIE} & \textbf{Ours} \\
\hline
High Res.              & \ding{55} & \ding{55} & \ding{55} & \ding{55} & \ding{51} & \ding{55} & \ding{55} & \ding{51} \\
Multimodal Labels      & \ding{55} & \ding{55} & \ding{51} & \ding{55} & \ding{51} & \ding{51} & \ding{55} & \ding{51} \\
Camera Pose            & \ding{55} & \ding{55} & \ding{55} & \ding{55} & \ding{51} & \ding{51} & \ding{55} & \ding{51} \\
Paired Images          & \ding{51} & \ding{51} & \ding{51} & \ding{51} & \ding{51} & \ding{51} & \ding{55} & \ding{51} \\
Multi Degrad.          & \ding{55} & \ding{55} & \ding{55} & \ding{51} & \ding{55} & \ding{55} & \ding{51} & \ding{51} \\
Physically Plausible   & \ding{55} & \ding{55} & \ding{55} & \ding{55} & \ding{51} & \ding{51} & \ding{51} & \ding{51} \\
\hline
\end{tabular}
\end{table*}

\subsection{光学多退化类型构建}

为模拟水下图像在不同水体条件下的真实退化特性，我们基于体积渲染模型构建了可调控的物理光传播机制，通过参数化设置实现多类型、多等级的退化图像生成。具体而言，我们在渲染过程中显式建模了光在水体传播过程中的吸收、前向散射和后向散射等核心物理机制，并在吸收系数 $\sigma_a(\lambda)$ 与散射系数 $\sigma_s(\lambda)$ 中引入波长依赖性变化，结合路径追踪（Path Tracing）算法与 Henyey–Greenstein 散射函数对光照进行物理一致性模拟，从而确保图像中的色偏、模糊和对比度下降等退化现象与真实海水环境保持高度一致。

具体地，光吸收表现为不同波长光线随传播距离呈指数衰减，我们通过调节不同波长通道（红、绿、蓝）的吸收系数$\sigma_a(\lambda)$来控制光谱特性，从而形成蓝绿色偏（增强红光吸收）和黄绿色偏（增强蓝绿光吸收）两种典型水下偏色方向。散射过程则由前向散射与后向散射组成，散射的强度与效果通过 Henyey–Greenstein 散射模型中的各向异性因子$g$与悬浮颗粒浓度（density）共同控制：较大的$g$值（如$0.6 \sim 0.9$）造成显著的前向散射，使图像边缘细节模糊；而后向散射产生的高强度背散射光形成了对比度降低、亮度漂白的视觉效果。上述光吸收与散射过程均与像素位置到摄像机的距离$d(x)$以及颗粒浓度呈强耦合关系，其退化效应的总体数学表达式可表示为经典的水下成像模型：
\begin{equation}
    I(x) = J(x)\cdot t(x) + B \cdot (1 - t(x)),
    \quad t(x) = \exp\big(-(\sigma_a+\sigma_s)\cdot d(x)\big)
\end{equation}
其中$I(x)$为最终的退化图像，$J(x)$为原始无退化图像，$B$为水下环境背景光强度，$t(x)$表示透射率，$\sigma_a$与$\sigma_s$分别表示吸收系数与散射系数。
基于上述物理模型，我们进一步构建了蓝色偏、蓝绿色偏、绿色偏、黄绿色偏、黄色偏五大类偏色方向，每类偏色方向下设置三种不同等级的浑浊程度，通过精确调节$\sigma_a(\lambda)$、$\sigma_s(\lambda)$、$g$值以及颗粒浓度，实现共十五种不同退化类型与等级的图像生成。这种系统性的参数控制，使生成的退化图像在结构保持一致的前提下，仅在色彩、模糊程度与对比度上产生变化，如图~\ref{fig:ours_degradation_levels}所示。此外，我们还结合了动态曝光调节机制，以适应真实相机在水下环境中的动态光照响应特性，使生成的图像更加真实自然。与现有水下数据集（如UIEB、EUVP）中基于风格迁移或经验生成的退化图像相比，本数据集中图像的退化过程更严格遵循真实物理规律，退化类型更全面，图像质量更具一致性与物理合理性，如图~\ref{fig:compare_degradation_uiieb}所示。

综上所述，我们提出的光学多退化类型建模策略不仅极大增强了合成数据的真实性与泛化能力，也为连续退化空间下的图像增强、色彩恢复与物理建模提供了结构对齐、参数可控且具备高度物理一致性的高质量训练基础。

\begin{figure*}[htbp]
    \centering
    \includegraphics[width=0.96\linewidth]{figs/ours_degradation_levels.pdf}
    \caption{(a)本数据集中两类典型偏色方向（蓝绿色偏、黄绿色偏）下的三种不同浑浊等级示意图，展示了退化类型的多样性和连续性。(b)与现有主流水下数据集生成的退化图像对比（左：UIEB/EUVP数据集中的退化图；右：本数据集物理建模生成的退化图）。可以看出本数据集中图像的结构、色彩和模糊程度更符合真实物理规律。}
    \label{fig:ours_degradation_levels}
\end{figure*}

% \begin{figure}[htbp]
%     \centering
%     \includegraphics[width=0.96\linewidth]{figs/compare_degradation_uiieb.pdf}
%     \caption{与现有主流水下数据集生成的退化图像对比（左：UIEB/EUVP数据集中的退化图；右：本数据集物理建模生成的退化图）。可以看出本数据集中图像的结构、色彩和模糊程度更符合真实物理规律。}
%     \label{fig:compare_degradation_uiieb}
% \end{figure}


\subsection{多模态及标注}

在数据集构建过程中，我们充分利用合成渲染流程的可控性与精确性，为每张图像同步生成了一套丰富且物理一致的多模态标注，以支持结构感知任务与跨模态建模需求。所有标注均来源于 Blender 渲染引擎中实时可导出的几何与相机参数，具有高精度、像素对齐和无噪声等特点，避免了真实数据集中常见的传感器误差与对齐偏移问题。

具体而言，我们为每一张图像输出以下标注信息：（1）\textbf{深度图（Depth map）}由 Z-buffer 导出，记录每个像素点在相机坐标系下的真实物理深度 $d(x)$，可直接用于单目深度估计、深度引导增强等任务；（2）\textbf{法线图（Normal map）}基于几何面片法线渲染得到，反映物体表面朝向与微观结构信息，在形状感知、光照建模等任务中具有关键作用；（3）\textbf{相机位姿（Camera pose）}记录渲染时相机的六自由度外参（位置与朝向）与完整内参矩阵（包括焦距、主点位置、像素比等参数），为多视图三维重建、新视角合成、视觉里程计等几何任务提供精准监督。

图~\ref{fig:multimodal_examples} 展示了我们为单帧图像生成的多模态标注示例，包括原始 RGB 图像、像素对齐的深度图与法线图，以及所对应的相机轨迹可视化结果。可以观察到，各模态间具备高度结构一致性，适用于多模态监督与跨任务特征协同建模。

为了获得具有几何一致性的多视角图像样本，我们在数据构建过程中设计了系统性的相机轨迹布置策略：相机在 $x$-$y$ 平面上构建位置网格，并在每个平移点上绕 $z$ 轴进行多角度旋转，从而采集结构一致但视角变化的数据图像，如图~\ref{fig:camera_sampling} 所示。该策略确保了图像之间的几何变换可追踪性，并为三维重建、视角合成等下游任务提供数据基础。

综上所述，本数据集在单帧图像增强之外，引入了可扩展、多用途的多模态几何标注体系，形成跨感知层级的完整监督机制，使其不仅适用于图像增强任务，也为深度估计、几何建图、位姿追踪等多项下游任务提供了统一、标准化的数据基础。

% \begin{figure}[htbp]
%     \centering
%     \includegraphics[width=0.98\linewidth]{figs/multimodal_examples.pdf}
%     \caption{多模态标注示例。依次为 RGB 图像、对应深度图、法线图。各模态信息结构对齐、物理一致，均来源于渲染引擎直接输出。}
%     \label{fig:multimodal_examples}
% \end{figure}

\begin{figure*}[htbp]
    \centering
    \includegraphics[width=0.85\linewidth]{camera_sampling.pdf}
    \caption{(a)多模态标注示例。依次为 RGB 图像、对应深度图、法线图。各模态信息结构对齐、物理一致，均来源于渲染引擎直接输出。(b)相机采样策略示意图。相机沿 $x$ 与 $y$ 平面网格布置，并在每个采样点上执行绕 $z$ 轴的旋转采样，形成结构一致但视角多样的图像序列。}
    \label{fig:camera_sampling}
\end{figure*}

\begin{figure*}[htbp]
    \centering
    \hspace{-0.17\linewidth}
    \includegraphics[width=0.83\linewidth]{datasetlabel.pdf}
    \caption{(a)主流水下图像增强数据集样本分辨率对比。我们选取了一组结构相似的珊瑚礁区域，选取红框区域放大八倍，可以看出我们的数据集在经过多倍放大后依然真实清晰，对照组数据集在经过等倍数放大后已经模糊不清无法分辨信息。(b)主流水下图像增强数据集样本数量以及平均分辨率对比。}
    \label{fig:datasetlabel}
\end{figure*}



% % 图 1: 数据集样本数对比图
% \begin{figure*}[htbp]
%     \centering
%     \includegraphics[width=0.95\linewidth]{dataset_comparison.pdf}
%     \caption{主流水下图像增强数据集样本分辨率对比。我们选取了一组结构相似的珊瑚礁区域，选取红框区域放大八倍，可以看出我们的数据集在经过多倍放大后依然真实清晰，对照组数据集在经过等倍数放大后已经模糊不清无法分辨信息。}
%     \label{fig:dataset_comparison}
% \end{figure*}

\subsection{数据集规模与统计}

本数据集在规模、分布与内容多样性方面均达到当前水下图像增强领域的领先水平。我们共构建了 $N_s$ 个高保真三维场景，每个场景在多个相机位置与不同退化条件下生成图像对，最终生成 RGB 图像、深度图与法线图共计 $N$ 对，每对图像之间具备严格的一一对应关系与像素级结构对齐。所有图像均采用 $3840 \times 2160$ 的 $4\text{K}$ 分辨率渲染，保持高细节还原度与真实物理一致性。图~\ref{fig:datasetlabel} 展示了我们数据集在图像分辨率及规模上的横向对比，包含了主流的水下图像增强数据集如 UIEB、LSUI等。可以看出，我们在保持高分辨率与多模态标注的前提下，实现了数量级上的跨越，为训练高容量模型提供了广覆盖的数据支持。



表~\ref{tab:dataset_comparison_vertical} 总结并对比了当前主流数据集的特征覆盖情况，包括高分辨率（High Res.）、多模态标签、相机位姿、结构对齐成对图像、多等级退化、多物理一致性等。
整体采样策略上，我们在 $x$-$y$ 平面构建 $M_x \times M_y$ 的二维位置网格，并在每个采样点上执行 $N_\text{yaw}$ 次绕 $z$ 轴旋转采样。每个视角对应图像都同步输出 RGB、深度、法线与位姿等完整信息，确保样本间几何变换可追踪，且具备一致结构基准，适用于多视角增强建模、几何一致性训练与跨帧监督机制构建。

综上所述，本数据集在数量规模、图像质量、多模态标签、几何结构保持以及退化配置丰富性等方面构成了当前水下图像增强任务中最具系统性与适应性的基准数据，为提升算法泛化能力与跨任务迁移效果提供了坚实的数据支撑。









\section{Proposed Method}

\subsection{Overall Framework（总框架与数据流）}
本文提出了一种面向多退化水下场景的端到端深度预测与图像增强网络，整体结构如图1所示。模型以\textbf{高效Transformer主干}为骨架，结合双分支协同重建和物理一致性引导，旨在同时实现高精度深度估计和退化图像的高质量恢复。整个网络采用\textbf{双次前向流程}，分别完成深度推理与增强重建两大子任务，彼此间实现紧密耦合与协同优化。

具体而言，输入退化水下图像 $I_{raw}$ 首先经过\textbf{浅层特征提取器（SFE）}，获得低级纹理与色彩特征。随后，特征送入\textbf{编码器模块}，多层Restormer结构提取逐级语义表征，并通过\textbf{跨模态双向交叉注意力机制}实现深度特征与RGB特征的相互引导与补充。首次前向阶段，以深度分支为主，利用解码器对多尺度编码特征进行上采与融合，输出连续深度估计 $\hat{D}$ 及多层深度特征。

在第二次前向中，模型以先验深度特征为辅助，对输入图像再次编码，在解码端设有\textbf{双物理一致性解码头}：分别针对去模糊与颜色校正进行分支设计。去模糊分支结合物理成像模型中的点扩散函数（PSF），实现深度相关的自适应模糊恢复；颜色校正分支基于Beer-Lambert衰减模型，对不同波段色彩衰减进行像素级建模。二者输出特征经融合后生成最终增强图像 $\hat{I}$。

最终，所有任务分支的输出经过残差融合与多任务头处理，端到端生成高质量增强图像与连续深度图，实现物理一致性约束下的多目标协同重建。网络训练过程中采用多项损失联合优化，利用不确定性加权机制自适应平衡多任务目标，显著提升了模型的泛化能力与训练稳定性。
\begin{figure*}[htbp]
  \centering
  \includegraphics[width=0.95\linewidth]{dualpipeline.pdf}
  \caption{dualpipeline}
  \label{fig:dualpipeline}
\end{figure*}
\begin{figure}[htbp]
  \centering
  \includegraphics[width=1\linewidth, trim=0cm 0cm 2cm 0cm, clip]{Bidirection_crossattention.pdf}
  \caption{Bidirection_crossattention}
  \label{fig:Bidirection_crossattention.pdf}
\end{figure}

\subsection{基于深度先验的双前向管线}
为了有效建模水下图像的复杂多退化过程及其与场景几何的高度耦合关系，本文提出了一种基于深度先验的双前向管线（Dual-Forward Pipeline with Depth Prior）。如图~\ref{fig:dualpipeline}所示该管线在一次完整推理中依次执行两次主干编码与分支解码，分别实现深度结构的建模与深度引导下的图像增强。



具体而言，输入退化水下图像 $\mathbf{I}_{\text{in}} \in \mathbb{R}^{H \times W \times 3}$ 首先经浅层特征提取器，依次通过两层 $3 \times 3$ 卷积、归一化与GELU激活，映射为统一通道数的低级特征 $\mathbf{F}_0 \in \mathbb{R}^{H \times W \times C}$。该特征作为主干编码器的输入，进入多层堆叠的Restormer块，每一层的输出表示为
\[
\mathbf{F}_l = \mathcal{B}_l(\mathbf{F}_{l-1}),
\]
其中 $\mathcal{B}_l$ 表示第 $l$ 层Restormer块，包括多头自注意力、前馈子网络及归一化等模块。主干编码器在两次前向中均保持参数共享，所有层的输出特征均被保存，用于后续分支解码与跨模态交互。

在第一次前向（Pass-1）阶段，编码器仅以输入特征 $\mathbf{F}_0$ 为流，依次输出多尺度特征 $\{\mathbf{F}_l^{(1)}\}_{l=1}^L$。该过程可递归表达为
\[
\mathbf{F}_l^{(1)} = \mathcal{B}_l(\mathbf{F}_{l-1}^{(1)}),\quad l=1,\ldots,L
\]
所有中间特征 $\{\mathbf{F}_l^{d}\}$ 及最终编码结果流向深度分支解码器，通过多层上采样与跳跃连接,编码器对称的多层Restomer模块及r2d交叉注意力机制重建出高分辨率连续深度图 $\hat{\mathbf{D}}$。所有深度中间特征在本阶段被缓存，作为后续增强分支的信息源。


在第二次前向（Pass-2）阶段，主干编码器输入不仅包含当前输入图像的初始特征 $\mathbf{F}_0$，还引入第一次前向阶段得到的多尺度深度特征 $\{\mathbf{F}_l^{d}\}$ 作为结构先验。编码器每一层均引入d2r交叉注意力机制，动态融合RGB流和深度流的信息。具体地，第 $l$ 层的RGB特征与深度特征分别表示为 $\mathbf{F}_l^{r}$ 和 $\mathbf{F}_l^{d}$，双向跨模态交叉注意力如图~\ref{fig:Bidirection_crossattention}交互更新过程如下：
\[
\begin{aligned}
\mathbf{F}_l^{r*} &= \mathrm{MultiHeadAttn}_r\left(\mathbf{Q}_r, \mathbf{K}_d, \mathbf{V}_d\right), \\
\mathbf{F}_l^{d*} &= \mathrm{MultiHeadAttn}_d\left(\mathbf{Q}_d, \mathbf{K}_r, \mathbf{V}_r\right), \\
\mathbf{F}_l^{r,\text{out}} &= \mathbf{F}_l^{r} + \mathbf{F}_l^{r*}, \\
\mathbf{F}_l^{d,\text{out}} &= \mathbf{F}_l^{d} + \mathbf{F}_l^{d*},
\end{aligned}
\]
其中 $\mathbf{Q}_r = W_q^r \mathbf{F}_l^{r}$, $\mathbf{K}_d = W_k^d \mathbf{F}_l^{d}$, $\mathbf{V}_d = W_v^d \mathbf{F}_l^{d}$ 等为线性投影得到的Query、Key、Value。反向方向同理。每层还可引入门控因子 $\gamma_l$ 对交互强度自适应调节：
\[
\mathbf{F}_l^{r,\text{out}} = \mathbf{F}_l^{r} + \gamma_l \cdot \mathbf{F}_l^{r*}
\]
门控参数随网络训练自动学习。融合后的多尺度特征被统一输入至增强分支解码器，该解码器内部采用双分支协同架构，分别针对水下图像的去模糊（Deblurring）与颜色校正（Color Correction）任务进行解码。

\subsection{双物理一致性解码与端到端重建}

\subsubsection{DPSF Deblurring Decoder}
针对水下成像中因多路径散射、颗粒介质和成像深度导致的空间相关模糊，本文提出了\textbf{深度引导点扩散函数去模糊解码器（Depth-guided PSF Deblurring Decoder, 简称DPSFD Decoder）}。该分支以深度信息为先验，充分结合点扩散函数（PSF）成像物理模型，在每一级解码过程中引入可学习的深度感知调制机制，实现了空间可变的自适应去模糊。

在结构上，DPSFD Decoder采用多级上采样与密集跳跃连接：每一级解码特征 $\mathbf{F}_d^{(i)}$，通过PixelShuffle与$1\times1$卷积进行空间分辨率恢复，并与编码器同层的多尺度RGB特征集成。融合后的特征张量 $\mathbf{Z}_d^{(i)}$，不仅携带了全局-局部的图像结构信息，更为后续物理建模提供了丰富的上下文表达。

在每一级解码中，DPSF分支嵌入了\textbf{深度引导PSF物理调制层（PSF-Physics Modulation Layer, PSF-PML）}。该层以归一化深度图 $D_{\mathrm{norm}}^{(i)}$ 与可学习模糊尺度参数 $\alpha$ 为条件，实现对特征流的像素级自适应调制：
\[
\sigma^{(i)} = \alpha \cdot D_{\mathrm{norm}}^{(i)}
\]
\[
\delta^{(i)} = \mathrm{MLP}\left(\mathrm{ReLU}\left(\mathbf{W}_1 \sigma^{(i)}\right)\right)
\]
\[
\mathbf{Z}_d^{(i)'} = \mathbf{Z}_d^{(i)} + \delta^{(i)}
\]
上述操作等效于在物理空间对理想图像 $J$ 进行空间变模糊的卷积过程：
\[
I_{\mathrm{obs}}(x, y) = \int\int J(u, v) \cdot h(x-u, y-v; \sigma(x, y))\, du\, dv
\]
其中 $h(\cdot,;\sigma)$ 为PSF核（常为高斯核），$\sigma$ 由深度自适应调节。不同于传统静态卷积，PSF-PML以深度场引导，动态模拟每个像素点处模糊核的变化，实现了空间非均质的模糊自适应去除。

融合物理感知调制后的特征$\mathbf{Z}_d^{(i)'}$，经过多层Restormer块进一步进行多头自注意力建模和门控前馈增强，提升对模糊纹理与结构信息的全局建模能力。最后，顶层特征通过卷积头映射输出高频去模糊残差$\mathbf{R}_D$：
\[
\mathbf{R}_D = f_{\mathrm{recon}}^{\mathrm{deblur}}\left(\mathbf{Z}_d^{(L)}\right)
\]
这一设计不仅保留了退化区域的结构纹理信息，更通过深度先验实现对空间变模糊的精准建模，为后续颜色恢复分支和全局融合提供了高质量的基础。

\subsubsection{DBL Color Correction Decoder}
针对水下环境中光谱衰减、散射等复杂光学效应造成的非均匀颜色失真，本文提出了\textbf{深度感知Beer-Lambert颜色恢复解码器（Depth-Aware Beer-Lambert Color Correction Decoder, 简称DBLCC Decoder）}。该分支以水下成像物理为核心，将深度信息与输入颜色紧密耦合，逐层引入\textbf{Beer-Lambert物理调制层（Beer-Lambert Physics Modulation Layer, BL-PML）}，对特征空间进行物理一致性修正。

BL-PML的设计严格遵循经典水下成像物理模型\cite{ref1}，即：
\[
I_c(x,y) = J_c(x,y) \cdot t_c(x,y) + B_c \cdot (1 - t_c(x,y))
\]
其中，$I_c(x,y)$为观测到的颜色通道退化像素，$J_c(x,y)$为无退化真实场景，$t_c(x,y)$为透射率（由深度$D(x,y)$和可学习消光系数$\beta_c$控制），$B_c$为背景光强度。透射率和后向散射项分别为：
\[
t_c(x,y) = \exp\left(-\beta_c \cdot D(x,y)\right), \quad b_c(x,y) = 1 - t_c(x,y)
\]
在网络实现上，DBLC Decoder采用多尺度特征上采与编码器跳连机制。每一级解码特征$\mathbf{F}*c^{(i)}$，融合来自编码器的多尺度语义与空间信息后，与原始输入图像$ I_{\mathrm{raw}} $及预测深度图$ D $共同送入BL-PML，实现像素级物理调制：
\[
\mathbf{F}_c^{(i)'} = \mathrm{BL}\text{-}\mathrm{PML}\left(\mathbf{F}_c^{(i)}, I_{\mathrm{raw}}^{(i)}, D^{(i)}, \beta_c, B_c\right)
\]
BL-PML在特征空间内以
\[
\mathbf{m}(x,y) = \left[t_c(x,y),\ b_c(x,y),\ 2I_{\mathrm{raw}}(x,y) - 1\right] \in \mathbb{R}^9
\]
为输入，经过小型卷积神经网络提取，映射为每层特征的仿射调制系数$(\gamma, \delta)$：
\[
[\gamma, \delta] = \mathrm{Conv}_{1\times1}\left(\mathrm{ReLU}(\mathrm{Conv}_{1\times1}(\mathbf{m}))\right)
\]
\[
\mathbf{F}_c^{(i)'} = \gamma \cdot \mathbf{F}_c^{(i)} + \delta
\]
上述物理一致性调制使解码器特征自适应响应于场景深度、局部光照和退化强度，实现了分布连续可调的像素级校正。紧随其后，多层\textbf{膨胀卷积-通道注意力模块（DACBlock）}进一步增强特征的多尺度上下文理解与通道自适应，充分捕获复杂水下场景中的色彩与结构信息。膨胀卷积凭借其扩展感受野能力，有效整合全局与细节的多尺度上下文，提升了网络对大范围颜色偏移与细粒度异常的建模能力。

最终，顶层解码特征经卷积重建输出颜色校正低频残差：
\[
\mathbf{R}_C = f_{\mathrm{recon}}^{\mathrm{color}}\left(\mathbf{F}_c^{(L)}\right)
\]
通过端到端训练，BL Color Decoder能够自适应学习深度-颜色物理规律下的光谱衰减与散射特性，实现高精度、物理一致性的颜色恢复，并极大提升了水下图像的可视化质量与后续处理的鲁棒性。

\subsubsection{分支融合与重建}
两条物理一致性分支输出的去模糊高频残差 $ \mathbf{R}_D $ 和颜色校正低频残差 $ \mathbf{R}*C $，与原始退化图像 $I_{\mathrm{raw}}$ 通过残差融合头进行端到端重建：
\[
\hat{I} = \phi\left(I_{\mathrm{raw}} + \mathbf{R}_D + \mathbf{R}_C\right)
\]
其中 $\phi(\cdot)$ 为逐元素tanh激活，保证输出图像范围与输入一致。



\subsection{损失函数与优化目标}
为实现高保真、结构感知强、色彩一致性高的水下图像增强效果，本文设计了一个多目标联合优化框架，涵盖图像重建、深度监督、物理及多模态一致性约束。最终的训练目标定义为：
\[
\mathcal{L}_{\text{total}} = \mathcal{L}_{\text{recon}} + \lambda_{\text{phys}} \mathcal{L}_{\text{phys}} + \lambda_{\text{distill}} \mathcal{L}_{\text{distill}} + \lambda_{\text{attn}} \mathcal{L}_{\text{attn}}
\]

\subsubsection{图像重建损失}
模型输出增强图像 $ \hat{I} = \tanh(I + \mathrm{res}_d + \mathrm{res}_c)$，其中结构残差与色彩残差分别由双分支生成。图像重建损失融合了像素、感知与结构信息，具体包括：
\[
\mathcal{L}_{\text{recon}} = \lambda_1 \|\hat{I} - I_{\text{gt}}\|_1 + \lambda_2\,\mathrm{SSIM} + \lambda_3\,\mathcal{L}_{\text{perc}} + \lambda_4\,\mathcal{L}_{\text{fft}} + \lambda_5\,\mathcal{L}_{\text{grad}}
\]
该设计在确保重建准确性的同时，提升图像感知质量与边缘清晰度。

\subsubsection{深度预测损失}

深度预测损失采用多项损失联合，定义如下：
\begin{equation}
\mathcal{L}_{\text{depth}} = \lambda_{d,1} \|\hat{D} - D_{\text{gt}}\|_1
+ \lambda_{d,2} \mathcal{L}_{\text{smooth}}
+ \lambda_{d,3} \mathcal{L}_{\text{edge}}
+ \lambda_{d,4} \mathcal{L}_{\text{ssim}}
\end{equation}
其中，分别对应像素级L1损失、深度平滑损失、结构边缘损失和结构相似性损失，用于提升全局深度精度、空间连续性、结构边界表现和整体结构一致性。各项权重$\lambda_{d,i}$采用不确定性自适应机制进行联合优化。




\subsubsection{跨模态协同损失}
为进一步强化RGB与深度模态间的结构对齐，我们引入交叉注意力一致性损失，用于约束Depth→RGB与RGB→Depth两个方向注意力图的对偶对称性：
\[
\mathcal{L}_{\text{attn}} = \sum_l \left\|A_{d2r}^{l} - (A_{r2d}^{l})^\top \right\|_1
\]
该损失有助于稳定跨模态交互过程，提升了最终结构恢复与物理一致性建模的效果。

\subsubsection{多退化一致性损失}

为抑制同一场景在不同退化条件下的输出波动，我们设计了基于深度置信度的像素方差约束。具体地，针对每个像素 $(x, y)$，我们计算其在 $N$ 种退化下增强结果的方差，并以预测深度为权重进行加权，得到一致性损失项：
% (2) 一致性损失
\begin{align}
\mathcal{L}_{\mathrm{var}}
&= \frac{1}{BHW} \sum_{b=1}^B \sum_{x,y}
    w_d^{(b)}(x,y)\cdot
    \mathrm{Var}\bigl(\{E^{(i)}_b(x,y)\}\bigr)
\end{align}

\begin{align}
w_d^{(b)}(x,y) &= \exp\bigl(-k\,\bar D_b(x,y)\bigr)
\end{align}

\begin{align}
\bar D_b(x,y) &= \frac{1}{N} \sum_{i=1}^N D^{(i)}_b(x,y)
\end{align}


其中，$w_d$ 表示深度置信权重，$k$ 为控制衰减幅度的常数，$\bar D$ 为多退化预测深度的均值。该项鼓励模型在浅水区域（低深度）保持更强的输出一致性，从而提升对复杂退化条件的鲁棒性。


此外，我们引入两项补充性一致性约束：其一，采用像素级 L1 损失对齐不同退化下的 RGB 增强结果，以抑制由退化差异引起的局部颜色漂移与结构偏差；其二，引入 L2 损失约束多退化下的深度预测结果保持一致，旨在提升深度估计对输入退化变化的鲁棒性。这两项设计共同促进模型在颜色还原与几何理解两个维度上实现跨退化一致的预测表现。


将上述三项线性组合，构成最终的多退化一致性损失函数：
\[
\mathcal{L}_{\mathrm{CMCL}} =
\lambda_{\mathrm{var}}\,\mathcal{L}_{\mathrm{var}} +
\lambda_{\mathrm{rgb}}\,\mathcal{L}_{\mathrm{rgb}} +
\lambda_{\mathrm{depth}}\,\mathcal{L}_{\mathrm{depth}}.
\]

\subsubsection{多任务不确定性加权机制}
为了在多任务联合训练中动态平衡各目标的重要性，本文进一步引入基于任务不确定性建模的损失加权策略：
\[
\mathcal{L}_{\text{total}} = \sum_k \left(\frac{1}{2\sigma_k^2} \mathcal{L}_k + \log \sigma_k\right)
\]
其中 $\sigma_k$ 为每个任务的可学习不确定性参数。该机制可有效缓解人工调参困难，提升多任务学习过程的稳定性与性能。

\section{Experiments}
在本节中，我们首先在现有主流水下图像增强数据集及所提出的 UBB 数据集上，对网络模型进行系统性的训练与评估，以全面分析数据集对模型性能的影响。此外，我们还将所提方法与当前最先进的单目深度估计算法进行了对比。最后，通过一系列消融实验，进一步验证了 Trident-Former 各组成模块的有效性与设计合理性。

\subsection{Experiment Setting}

\textbf{Datasets.} 本文主要在多个主流水下图像增强与深度估计数据集上进行实验。我们提出的 UBB 数据集被随机划分为训练集和测试集，分别用于模型训练与性能评估。除此之外，实验还使用了 LSUI（4500 对图像）、EUVP（2185 对图像）和 UIEB（800 对图像）作为额外训练数据，以提升模型的泛化能力和多场景适应性。测试集分为以下三类：

\begin{itemize}
\item 有参考测试集（Reference-based）：包括 testUBB、TestU90（由 UIEB 中 90 张具备参考图像的样本组成）、TestL400（由LSUI中400有参考图像组成）。
\item 无参考测试集（Reference-free）：包括 testUBB、Test-U60（由 UIEB 中 60 张无参考图像组成）、SQUID（共 16 张图像）。
% \item 深度评测集（Depth Evaluation）：涵盖 testUBB、Sea-thru D3、Sea-thru D5 和 SQUID 四个子集，分别包含 43、48 和 57 张具有真实深度标签的图像。
\end{itemize}

\textbf{Compared Methods.}
在图像增强实验中，选取了代表性的物理模型方法和多种最新深度学习方法进行对比。在深度估计实验中，分引入了通用的预训练深度模型（如 MegaDepth、MiDaS、DPT）以全面评估增强和深度估计任务的性能上限与现有水平。

\textbf{Evaluation Metrics.}
实验采用多种客观评价指标，覆盖图像增强与深度估计两个方向：

\begin{itemize}
\item \textbf{图像增强指标：} 包括全参考指标和无参考指标两类。全参考指标包括峰值信噪比（PSNR）、结构相似性（SSIM）、感知相似度（LPIPS）以及色彩差异度量（CIEDE2000）。其中，PSNR 衡量增强图像与参考图像的像素级误差，SSIM 强调图像结构的一致性，LPIPS 衡量感知空间下的视觉相似性，CIEDE2000 则反映色彩恢复的准确性。无参考指标包括水下图像质量评价指数（UCIQE）、水下图像质量度量（UIQM）以及自然图像质量评价指标（NIQE），分别从色彩、对比度、锐度与自然度等方面量化增强效果。

\item \textbf{深度估计指标：} 包括精度指标与误差指标两大类。精度指标采用尺度阈值准确率 $\delta < 1.05$、$\delta < 1.05^2$、$\delta < 1.05^3$，分别统计预测深度与真实深度在各阈值下的像素比例。误差指标包括绝对相对误差（AbsRel）、平方相对误差（Sq.Rel）、均方根误差（RMSE）、对数均方根误差（RMSE-log）以及尺度不变对数误差（SILog），用于全面衡量深度估计结果的准确性与稳健性。
\end{itemize}


所有实验均在同一硬件环境和训练设置下进行，具体实现细节见附录 A。

\subsection{Dataset Evaluation}
为分析不同的训练数据对模型泛化性能的影响，我们选取了多种主流水下图像增强方法（包括waternet、ucolor、PUIE-Net、NU2Net、U-Shape Transformer）分别在UBB、LSUI、EUVP、和UIEB水下图像增强数据集上独立训练，并在统一测试协议下于多个标准测试集（如 Test-U90、TestL400、Test-U60、SQUID）进行横向评测。
\begin{table*}[htbp]
\centering
\scriptsize
\caption{Cross‐dataset comparison on \textbf{reference‐based} test sets.
Best in each column is \textbf{bold}.}
\label{tab:ref_all_metrics}
\setlength{\tabcolsep}{4pt}
\begin{tabular}{|c|c|
    c c c c|  % testUBB: PSNR, SSIM, CIEDE2000, LPIPS
    c c c c|  % TestU90
    c c c c|  % TestL400
|}
\hline
\multirow{\textbf{Model}}
& \multirow{\textbf{Train}}
& \multicolumn{4}{c|}{\textbf{testUBB}}
& \multicolumn{4}{c|}{\textbf{TestU90}}
& \multicolumn{4}{c|}{\textbf{TestL400}} \\
\cline{3-14}
& & PSNR↑ & SSIM↑ & CIEDE↓ & LPIPS↓
    & PSNR↑ & SSIM↑ & CIEDE↓ & LPIPS↓
    & PSNR↑ & SSIM↑ & CIEDE↓ & LPIPS↓ \\ \hline

\multirow{WaterNet}
  & UBB  & 〈〉 & 〈〉 & 〈〉 & 〈〉
         & 〈〉 & 〈〉 & 〈〉 & 〈〉
         & 〈〉 & 〈〉 & 〈〉 & 〈〉 \\
  & LSUI & 〈〉 & 〈〉 & 〈〉 & 〈〉
         & 〈〉 & 〈〉 & 〈〉 & 〈〉
         & 〈〉 & 〈〉 & 〈〉 & 〈〉 \\
  & EUVP & 〈〉 & 〈〉 & 〈〉 & 〈〉
         & 〈〉 & 〈〉 & 〈〉 & 〈〉
         & 〈〉 & 〈〉 & 〈〉 & 〈〉 \\
  & UIEB & 〈〉 & 〈〉 & 〈〉 & 〈〉
         & 〈〉 & 〈〉 & 〈〉 & 〈〉
         & 〈〉 & 〈〉 & 〈〉 & 〈〉 \\ \hline

\multirow{Ucolor}
  & UBB  & 〈〉 & 〈〉 & 〈〉 & 〈〉
         & 〈〉 & 〈〉 & 〈〉 & 〈〉
         & 〈〉 & 〈〉 & 〈〉 & 〈〉 \\
  & LSUI & 〈〉 & 〈〉 & 〈〉 & 〈〉
         & 〈〉 & 〈〉 & 〈〉 & 〈〉
         & 〈〉 & 〈〉 & 〈〉 & 〈〉 \\
  & EUVP & 〈〉 & 〈〉 & 〈〉 & 〈〉
         & 〈〉 & 〈〉 & 〈〉 & 〈〉
         & 〈〉 & 〈〉 & 〈〉 & 〈〉 \\
  & UIEB & 〈〉 & 〈〉 & 〈〉 & 〈〉
         & 〈〉 & 〈〉 & 〈〉 & 〈〉
         & 〈〉 & 〈〉 & 〈〉 & 〈〉 \\ \hline

\multirow{PUIE-Net}
  & UBB  & 〈〉 & 〈〉 & 〈〉 & 〈〉
         & 〈〉 & 〈〉 & 〈〉 & 〈〉
         & 〈〉 & 〈〉 & 〈〉 & 〈〉 \\
  & LSUI & 〈〉 & 〈〉 & 〈〉 & 〈〉
         & 〈〉 & 〈〉 & 〈〉 & 〈〉
         & 〈〉 & 〈〉 & 〈〉 & 〈〉 \\
  & EUVP & 〈〉 & 〈〉 & 〈〉 & 〈〉
         & 〈〉 & 〈〉 & 〈〉 & 〈〉
         & 〈〉 & 〈〉 & 〈〉 & 〈〉 \\
  & UIEB & 〈〉 & 〈〉 & 〈〉 & 〈〉
         & 〈〉 & 〈〉 & 〈〉 & 〈〉
         & 〈〉 & 〈〉 & 〈〉 & 〈〉 \\ \hline

\multirow{NU2Net}
  & UBB  & 〈〉 & 〈〉 & 〈〉 & 〈〉
         & 〈〉 & 〈〉 & 〈〉 & 〈〉
         & 〈〉 & 〈〉 & 〈〉 & 〈〉 \\
  & LSUI & 〈〉 & 〈〉 & 〈〉 & 〈〉
         & 〈〉 & 〈〉 & 〈〉 & 〈〉
         & 〈〉 & 〈〉 & 〈〉 & 〈〉 \\
  & EUVP & 〈〉 & 〈〉 & 〈〉 & 〈〉
         & 〈〉 & 〈〉 & 〈〉 & 〈〉
         & 〈〉 & 〈〉 & 〈〉 & 〈〉 \\
  & UIEB & 〈〉 & 〈〉 & 〈〉 & 〈〉
         & 〈〉 & 〈〉 & 〈〉 & 〈〉
         & 〈〉 & 〈〉 & 〈〉 & 〈〉 \\ \hline

\multirow{U-Shape Transformer}
  & UBB  & 〈〉 & 〈〉 & 〈〉 & 〈〉
         & 〈〉 & 〈〉 & 〈〉 & 〈〉
         & 〈〉 & 〈〉 & 〈〉 & 〈〉 \\
  & LSUI & 〈〉 & 〈〉 & 〈〉 & 〈〉
         & 〈〉 & 〈〉 & 〈〉 & 〈〉
         & 〈〉 & 〈〉 & 〈〉 & 〈〉 \\
  & EUVP & 〈〉 & 〈〉 & 〈〉 & 〈〉
         & 〈〉 & 〈〉 & 〈〉 & 〈〉
         & 〈〉 & 〈〉 & 〈〉 & 〈〉 \\
  & UIEB & 〈〉 & 〈〉 & 〈〉 & 〈〉
         & 〈〉 & 〈〉 & 〈〉 & 〈〉
         & 〈〉 & 〈〉 & 〈〉 & 〈〉 \\ \hline

\multirow{\textbf{Trident-Former}}
  & UBB  & \textbf{〈〉} & \textbf{〈〉} & \textbf{〈〉} & \textbf{〈〉}
         & \textbf{〈〉} & \textbf{〈〉} & \textbf{〈〉} & \textbf{〈〉}
         & \textbf{〈〉} & \textbf{〈〉} & \textbf{〈〉} & \textbf{〈〉} \\
  & LSUI & 〈〉 & 〈〉 & 〈〉 & 〈〉
         & 〈〉 & 〈〉 & 〈〉 & 〈〉
         & 〈〉 & 〈〉 & 〈〉 & 〈〉 \\
  & EUVP & 〈〉 & 〈〉 & 〈〉 & 〈〉
         & 〈〉 & 〈〉 & 〈〉 & 〈〉
         & 〈〉 & 〈〉 & 〈〉 & 〈〉 \\
  & UIEB & 〈〉 & 〈〉 & 〈〉 & 〈〉
         & 〈〉 & 〈〉 & 〈〉 & 〈〉
         & 〈〉 & 〈〉 & 〈〉 & 〈〉 \\ \hline
\end{tabular}
\end{table*}


\begin{table*}[htbp]
\centering
\scriptsize
\caption{Cross‐dataset comparison on \textbf{reference‐free} test sets.
Best in each column is \textbf{bold}.}
\label{tab:noref_all}
\setlength{\tabcolsep}{8pt}
\begin{tabular}{|c|c|
    c c c|  % testUBB: UCIQE, UIQM, NIQE
    c c c|  % Test-U60
    c c c|  % SQUID
|}
\hline
\multirow{\textbf{Model}}
& \multirow{\textbf{Train}}
& \multicolumn{3}{c|}{\textbf{testUBB}}
& \multicolumn{3}{c|}{\textbf{Test-U60}}
& \multicolumn{3}{c|}{\textbf{SQUID}} \\
\cline{3-11}
& & UCIQE↑ & UIQM↑ & NIQE↓
    & UCIQE↑ & UIQM↑ & NIQE↓
    & UCIQE↑ & UIQM↑ & NIQE↓ \\ \hline

\multirow{WaterNet}
  & UBB  & 〈〉 & 〈〉 & 〈〉
         & 〈〉 & 〈〉 & 〈〉
         & 〈〉 & 〈〉 & 〈〉 \\
  & LSUI & 〈〉 & 〈〉 & 〈〉
         & 〈〉 & 〈〉 & 〈〉
         & 〈〉 & 〈〉 & 〈〉 \\
  & EUVP & 〈〉 & 〈〉 & 〈〉
         & 〈〉 & 〈〉 & 〈〉
         & 〈〉 & 〈〉 & 〈〉 \\
  & UIEB & 〈〉 & 〈〉 & 〈〉
         & 〈〉 & 〈〉 & 〈〉
         & 〈〉 & 〈〉 & 〈〉 \\ \hline

\multirow{Ucolor}
  & UBB  & 〈〉 & 〈〉 & 〈〉
         & 〈〉 & 〈〉 & 〈〉
         & 〈〉 & 〈〉 & 〈〉 \\
  & LSUI & 〈〉 & 〈〉 & 〈〉
         & 〈〉 & 〈〉 & 〈〉
         & 〈〉 & 〈〉 & 〈〉 \\
  & EUVP & 〈〉 & 〈〉 & 〈〉
         & 〈〉 & 〈〉 & 〈〉
         & 〈〉 & 〈〉 & 〈〉 \\
  & UIEB & 〈〉 & 〈〉 & 〈〉
         & 〈〉 & 〈〉 & 〈〉
         & 〈〉 & 〈〉 & 〈〉 \\ \hline

\multirow{PUIE-Net}
  & UBB  & 〈〉 & 〈〉 & 〈〉
         & 〈〉 & 〈〉 & 〈〉
         & 〈〉 & 〈〉 & 〈〉 \\
  & LSUI & 〈〉 & 〈〉 & 〈〉
         & 〈〉 & 〈〉 & 〈〉
         & 〈〉 & 〈〉 & 〈〉 \\
  & EUVP & 〈〉 & 〈〉 & 〈〉
         & 〈〉 & 〈〉 & 〈〉
         & 〈〉 & 〈〉 & 〈〉 \\
  & UIEB & 〈〉 & 〈〉 & 〈〉
         & 〈〉 & 〈〉 & 〈〉
         & 〈〉 & 〈〉 & 〈〉 \\ \hline

\multirow{NU2Net}
  & UBB  & 〈〉 & 〈〉 & 〈〉
         & 〈〉 & 〈〉 & 〈〉
         & 〈〉 & 〈〉 & 〈〉 \\
  & LSUI & 〈〉 & 〈〉 & 〈〉
         & 〈〉 & 〈〉 & 〈〉
         & 〈〉 & 〈〉 & 〈〉 \\
  & EUVP & 〈〉 & 〈〉 & 〈〉
         & 〈〉 & 〈〉 & 〈〉
         & 〈〉 & 〈〉 & 〈〉 \\
  & UIEB & 〈〉 & 〈〉 & 〈〉
         & 〈〉 & 〈〉 & 〈〉
         & 〈〉 & 〈〉 & 〈〉 \\ \hline

\multirow{U-Shape Transformer}
  & UBB  & 〈〉 & 〈〉 & 〈〉
         & 〈〉 & 〈〉 & 〈〉
         & 〈〉 & 〈〉 & 〈〉 \\
  & LSUI & 〈〉 & 〈〉 & 〈〉
         & 〈〉 & 〈〉 & 〈〉
         & 〈〉 & 〈〉 & 〈〉 \\
  & EUVP & 〈〉 & 〈〉 & 〈〉
         & 〈〉 & 〈〉 & 〈〉
         & 〈〉 & 〈〉 & 〈〉 \\
  & UIEB & 〈〉 & 〈〉 & 〈〉
         & 〈〉 & 〈〉 & 〈〉
         & 〈〉 & 〈〉 & 〈〉 \\ \hline

\multirow{\textbf{Trident-Former}}
  & UBB  & \textbf{〈〉} & \textbf{〈〉} & \textbf{〈〉}
         & \textbf{〈〉} & \textbf{〈〉} & \textbf{〈〉}
         & \textbf{〈〉} & \textbf{〈〉} & \textbf{〈〉} \\
  & LSUI & 〈〉 & 〈〉 & 〈〉
         & 〈〉 & 〈〉 & 〈〉
         & 〈〉 & 〈〉 & 〈〉 \\
  & EUVP & 〈〉 & 〈〉 & 〈〉
         & 〈〉 & 〈〉 & 〈〉
         & 〈〉 & 〈〉 & 〈〉 \\
  & UIEB & 〈〉 & 〈〉 & 〈〉
         & 〈〉 & 〈〉 & 〈〉
         & 〈〉 & 〈〉 & 〈〉 \\ \hline

\end{tabular}
\end{table*}

实验结果如表 X 所示。可以观察到，绝大多数增强模型在 UBB 数据集上训练后，均在全参考指标（PSNR、SSIM、CIEDE2000、LPIPS）与无参考指标（UCIQE、UIQM、NIQE）上取得最优或有竞争力的性能，尤其在跨数据集的泛化测试中（如 SQUID），优势更加显著。这充分证明了 UBB 数据集在数据多样性、物理一致性和高质量标注等方面的领先性。
图~\ref{fig:trident_multiset_sample} 展示了 Trident-Former 分别在多个训练数据集（如 UBB、LSUI、EUVP 等）上训练后所得到的图像增强结果，可作为本文数据评估部分的补充参考。从图中可以观察到，使用 UBB 数据集训练的模型在多数测试图像上展现出最佳的指标性能和较为自然的视觉观感；而在其他数据集上训练的模型则不同程度地出现偏色、过饱和或细节缺失等问题。这一现象进一步证明了 UBB 数据集因其符合水下光学成像物理规律、覆盖多样退化场景的特性,精准对齐的高分辨率多模态标注，使得模型在泛化能力上具备显著优势，能够更有效地适应复杂多变的水下环境。
我们进一步分析认为，UBB 的提升得益于：(1) 丰富的场景类型与真实物理过程建模，提升了网络的适应性和泛化能力；(2) 精准对齐的高分辨率多模态标注，显著增强了模型的学习能力；(3) 多退化等级和全覆盖采样策略，为各类增强方法提供了坚实的数据基础。综上，UBB 数据集为推动水下图像增强及相关视觉任务的发展提供了有力支撑。

\subsection{Network Architecture Evaluation}
前小节已经做过Trident-Former与多个先进水下图像增强模型在多个训练集、测试集以及参考指标下的数据对比，本小节将进一步引进多个具有代表性的水下图像增强算法并进一步评估Trident-Former在多个维度上的优势，所有测试均基于UBB数据集训练，在TestUBB，TestU90、TestL400上进行参考指标对比，在TestUBB、、Test-U60、SQUID数据集上进行无参考指标对比。
\begin{table*}[t]
\centering
\caption{十种方法在全参考测试集上的定量对比（PSNR$\uparrow$, SSIM$\uparrow$, LPIPS$\downarrow$, CIEDE2000$\downarrow$）}
\label{tab:full_ref_compact}
\setlength\tabcolsep{2pt}
\resizebox{0.95\textwidth}{!}{
\begin{tabular}{l|cccc|cccc|cccc}
\toprule
\multirow{方法}
 & \multicolumn{4}{c|}{testUBB}
 & \multicolumn{4}{c|}{TestU90}
 & \multicolumn{4}{c}{TestL400} \\
 & PSNR$\uparrow$ & SSIM$\uparrow$ & LPIPS$\downarrow$ & CIEDE$\downarrow$
 & PSNR$\uparrow$ & SSIM$\uparrow$ & LPIPS$\downarrow$ & CIEDE$\downarrow$
 & PSNR$\uparrow$ & SSIM$\uparrow$ & LPIPS$\downarrow$ & CIEDE$\downarrow$ \\
\midrule
Fusion         & –    & –    & –     & –      & –    & –    & –     & –      & –    & –    & –     & –      \\
SMBL           & –    & –    & –     & –      & –    & –    & –     & –      & –    & –    & –     & –      \\
MLLE           & –    & –    & –     & –      & –    & –    & –     & –      & –    & –    & –     & –      \\
UWCNN          & –    & –    & –     & –      & –    & –    & –     & –      & –    & –    & –     & –      \\
Water-Net      & –    & –    & –     & –      & –    & –    & –     & –      & –    & –    & –     & –      \\
Ucolor         & –    & –    & –     & –      & –    & –    & –     & –      & –    & –    & –     & –      \\
PUIE-Net       & –    & –    & –     & –      & –    & –    & –     & –      & –    & –    & –     & –      \\
TACL           & –    & –    & –     & –      & –    & –    & –     & –      & –    & –    & –     & –      \\
NU2Net         & –    & –    & –     & –      & –    & –    & –     & –      & –    & –    & –     & –      \\
U-shape Trans. & –    & –    & –     & –      & –    & –    & –     & –      & –    & –    & –     & –      \\
\bottomrule
\end{tabular}
}
\end{table*}


\begin{table*}[t]
\centering
\caption{十种方法在无参考测试集上的定量对比（UCIQE$\uparrow$, UIQM$\uparrow$, NIQE$\downarrow$）}
\label{tab:no_ref_comparison}
\setlength\tabcolsep{3pt}
\begin{tabular}{l|ccc|ccc|ccc}
\toprule
\multirow{方法}
 & \multicolumn{3}{c|}{testUBB}
 & \multicolumn{3}{c|}{Test-U60}
 & \multicolumn{3}{c}{SQUID} \\
 & UCIQE$\uparrow$ & UIQM$\uparrow$ & NIQE$\downarrow$
 & UCIQE$\uparrow$ & UIQM$\uparrow$ & NIQE$\downarrow$
 & UCIQE$\uparrow$ & UIQM$\uparrow$ & NIQE$\downarrow$ \\
\midrule
Fusion         & –    & –    & –    & –    & –    & –    & –    & –    & –    \\
SMBL           & –    & –    & –    & –    & –    & –    & –    & –    & –    \\
MLLE           & –    & –    & –    & –    & –    & –    & –    & –    & –    \\
UWCNN          & –    & –    & –    & –    & –    & –    & –    & –    & –    \\
Water-Net      & –    & –    & –    & –    & –    & –    & –    & –    & –    \\
Ucolor         & –    & –    & –    & –    & –    & –    & –    & –    & –    \\
PUIE-Net       & –    & –    & –    & –    & –    & –    & –    & –    & –    \\
TACL           & –    & –    & –    & –    & –    & –    & –    & –    & –    \\
NU2Net         & –    & –    & –    & –    & –    & –    & –    & –    & –    \\
U-shape Trans. & –    & –    & –    & –    & –    & –    & –    & –    & –    \\
\bottomrule
\end{tabular}
\end{table*}

接下来我们将对Trident-Former联合UBB数据集在深度估计上的性能与 MegaDepth、MiDaS、DPT进行对比如表*。可以看出Trident-Former在深度估计方面已经达到了优秀水准，通过多任务联合优化和双向交叉注意力机制，实现了图像增强与深度估计的协同学习，相比传统单任务方法在水下复杂退化场景下表现出更强的鲁棒性和更高的预测精度，特别是在深度边缘保持和物理一致性方面显著优于MiDaS等空中预训练模型。同时也体现了UBB数据集强大的泛化能力以及其高保真物理建模和多退化等级设计为复杂水下深度估计任务提供了更加丰富和可靠的监督信号，有效缓解了传统方法因训练数据域偏移导致的性能下降问题。
\begin{table*}[htbp]
  \centering
  \small % 调整字体为小号
  \setlength{\tabcolsep}{3pt} % 调整列间距
  \caption{Trident-Former（UBB 数据集）与 MegaDepth、MiDaS、DPT 的深度估计性能对比}
  \label{tab:depth_comparison}
  \begin{tabular}{lcccccccc}
    \toprule
    方法 & $\delta<1.05 \uparrow$ & $\delta<1.05^2 \uparrow$ & $\delta<1.05^3 \uparrow$ & AbsRel$\downarrow$ & Sq.Rel$\downarrow$ & RMSE$\downarrow$ & RMSE-log$\downarrow$ & SILog$\downarrow$ \\
    \midrule
    Trident-Former + UBB &   &   &   &   &   &   &   &   \\
    MegaDepth            &   &   &   &   &   &   &   &   \\
    MiDaS                &   &   &   &   &   &   &   &   \\
    DPT                  &   &   &   &   &   &   &   &   \\
    \bottomrule
  \end{tabular}
\end{table*}

\subsection{Ablation Studies}
（消融实验）
为验证各组模块的有效性，我们在UBB的有参考数据集以及无参考数据集上进行了系列消融实验。实验考虑了四个关键因素：depth guide，DPSF Deblurring Decoder，DBL Color Correction decoder，多退化一致性损失。
\begin{table*}[htbp]
  \centering
  \small
  \setlength{\tabcolsep}{5pt}
  \caption{各模块消融实验结果。左侧为UBB有参考测试集（PSNR、SSIM、LPIPS、CIEDE），右侧为无参考测试集（UCIQE、UIQM、NIQE）。}
  \label{tab:ablation}
  \begin{tabular}{lcccc|ccc}
    \toprule
    方法 & PSNR$\uparrow$ & SSIM$\uparrow$ & LPIPS$\downarrow$ & CIEDE$\downarrow$ & UCIQE$\uparrow$ & UIQM$\uparrow$ & NIQE$\downarrow$ \\
    \midrule
    Ours (Full)                 &       &       &       &       &       &       &       \\
    w/o Depth Guide             &       &       &       &       &       &       &       \\
    w/o DPSF Deblurring Decoder &       &       &       &       &       &       &       \\
    w/o DBL Color Correction    &       &       &       &       &       &       &       \\
    w/o Multi-degradation Loss  &       &       &       &       &       &       &       \\
    \bottomrule
  \end{tabular}
\end{table*}



% --- 参考文献 ---
\begin{thebibliography}{99}
\bibitem{ref1} Placeholder 1
\bibitem{ref2} Placeholder 2
\bibitem{ref3} Placeholder 3
\bibitem{ref4} Placeholder 4
\bibitem{ref5} Placeholder 5
\bibitem{ref6} Placeholder 6
\bibitem{ref7} Placeholder 7
\bibitem{ref8} Placeholder 8
\bibitem{ref9} Placeholder 9
\bibitem{ref16} Placeholder 16
\bibitem{ref17} Placeholder 17
\bibitem{ref18} Placeholder 18
\bibitem{ref19} Placeholder 19
\bibitem{ref21} Placeholder 21
\bibitem{ref23} Placeholder 23
\bibitem{ref24} Placeholder 24
\end{thebibliography}
% 为验证所提出的 Trident-Former 框架的有效性及 UBB 数据集所带来的性能提升，我们设计了以下四个方面的实验：

% 1. \textbf{跨数据集训练与泛化能力评估}：为系统分析不同训练数据对模型泛化性能的影响，我们分别在 UBB、LSUI、EUVP 和 UIEB 等主流水下图像增强数据集上单独训练 Trident-Former，并在所有数据集划分的测试集（包括 testUBB、Test-U90、Test-U60、SQUID 等）进行交叉测试。这一实验旨在定量衡量训练数据集对模型泛化能力的影响，突出 UBB 数据集在数据分布覆盖广度与退化模型泛化适应性方面的优势。

% 2. \textbf{主流方法对比实验}：为全面评估 Trident-Former 在水下图像增强任务中的性能优势，我们选取了当前最具代表性的多种方法进行横向对比，包括：

%    * 基于物理模型的方法：UIBLA、UDCP；
%    * 基于视觉先验的方法：Fusion、Retinex-based、RGHS；
%    * 基于深度学习的数据驱动方法：U-Shape Transformer、Ucolor、WaterNet、UIE-DAL、FUnIE-GAN。

%    我们在上述划分的各测试集上分别进行定量评估，所采用的评测指标包括全参考指标与无参考指标。全参考指标包括峰值信噪比（Peak Signal-to-Noise Ratio, PSNR）与结构相似性指数（Structural Similarity Index Measure, SSIM），前者衡量增强图像与参考图像间的像素级差异，后者则强调图像结构信息的一致性。无参考图像质量评价指标包括水下图像质量评价指数（Underwater Color Image Quality Evaluation, UCIQE）、水下图像质量度量（Underwater Image Quality Measure, UIQM）以及自然图像质量评价指标（Natural Image Quality Evaluator, NIQE）。UCIQE 和 UIQM 主要针对水下图像的颜色饱和度、亮度对比度和清晰度等特性进行量化，NIQE 则衡量增强后图像的自然感知质量。此外，我们还进行了主观视觉评分（Perception Score, PS），由多位受训观察者在1到5的范围内对增强结果进行评价，以辅助客观指标更全面地体现模型性能。

% 3. \textbf{深度估计任务迁移实验}：为进一步验证 Trident-Former 增强图像对下游计算机视觉任务的实际提升效果，我们将增强后的图像作为输入，分别使用 MiDaS、DPT 和 PUDE 等先进的单目深度估计模型进行评估。我们使用深度估计任务中的常用指标进行定量评测，包括绝对相对误差（Absolute Relative Error, AbsRel）、平方相对误差（Square Relative Error, Sq.Rel）、均方根误差（Root Mean Squared Error, RMSE）、对数空间均方根误差（RMSE log）和尺度不变对数误差（Scale-Invariant Logarithmic Error, SILog），以及准确度指标（\$\delta < 1.05\$、\$\delta < 1.05^2\$ 和 \$\delta < 1.05^3\$），以充分衡量图像增强质量对下游深度估计性能的影响。

% 4. \textbf{消融实验}：为量化分析 Trident-Former 各核心模块与 UBB 数据集的关键属性对最终性能的贡献，我们分别进行模块级与数据级消融实验。模块级消融主要关注物理一致性双分支解码器、跨模态注意力机制以及深度引导机制等关键设计；数据级消融则分别移除多模态标注信息、高分辨率特性、多等级退化控制等数据集关键属性，单独或组合进行训练评估，以明确各组成部分的贡献。

% 以上所有实验均采用统一的训练超参数与评测标准，训练与实现细节参见附录 A，以保证实验结果的公正性、可比性与可复现性。

\end{document} 